% =====================================================================
%  Namespaces Without a Name Server
%  L. G. Meredith, F1R3FLY.io
% =====================================================================
\documentclass[11pt]{article}

\usepackage[T1]{fontenc}
\usepackage[utf8]{inputenc}
\usepackage{amsmath,amssymb,amsthm}
\usepackage[dvipsnames]{xcolor}
\usepackage[margin=1.1in]{geometry}
\usepackage{mathrsfs}
\usepackage{enumitem}
\usepackage{booktabs}
\usepackage{listings}
\usepackage[expansion=false]{microtype}
\usepackage[hidelinks]{hyperref}

\theoremstyle{definition}
\newtheorem{definition}{Definition}[section]
\newtheorem{construction}[definition]{Construction}
\newtheorem{notation}[definition]{Notation}
\theoremstyle{plain}
\newtheorem{proposition}[definition]{Proposition}
\newtheorem{lemma}[definition]{Lemma}
\newtheorem{theorem}[definition]{Theorem}
\newtheorem{corollary}[definition]{Corollary}
\theoremstyle{remark}
\newtheorem{remark}[definition]{Remark}
\newtheorem{example}[definition]{Example}

% ---- notation ----
\newcommand{\rhoc}{\mathrm{rho}}
\newcommand{\quo}[1]{@#1}
\newcommand{\drp}[1]{{*}#1}
\newcommand{\inp}[3]{\mathsf{for}(\,#1 \leftarrow #2\,)#3}
\newcommand{\inpp}[5]{\mathsf{for}(\,#1 \leftarrow #2\ \&\ #3 \leftarrow #4\,)#5}
\newcommand{\out}[2]{#1!(#2)}
\newcommand{\enc}[1]{[\![#1]\!]}
\newcommand{\nequiv}{\equiv_{\mathcal{N}}}
\newcommand{\nil}{\mathbf{0}}
\newcommand{\Lft}{\mathsf{L}}
\newcommand{\Rgt}{\mathsf{R}}
\newcommand{\Frsh}{\mathsf{Fresh}_2}
\newcommand{\Dup}{\mathsf{D}}
\newcommand{\Ctx}{\mathscr{C}}
\newcommand{\Obs}{\mathcal{D}}
\newcommand{\pic}{\pi}
\newcommand{\bisim}{\approx}
\newcommand{\QD}{\mathrm{qd}}
\newcommand{\fn}{\mathrm{fn}}
\providecommand{\xLongrightarrow}[1]{\mathrel{\overset{#1}{\Longrightarrow}}}

\lstset{basicstyle=\ttfamily\small,columns=fullflexible,keepspaces=true,
        breaklines=true,frame=single,framesep=4pt,xleftmargin=6pt}

\title{\textbf{Namespaces Without a Name Server}\\[4pt]
\large Contraction, copying, and a fully abstract encoding of the
$\pi$-calculus in the rho calculus}
\author{L.\,G.\ Meredith\thanks{F1R3FLY.io.\ \texttt{lgreg.meredith@f1r3fly.io}}}
\date{August 2026}

\begin{document}
\maketitle

\begin{abstract}
\noindent
The rho calculus was designed so that its theory of names would be
\emph{internal}: names are quoted processes, there is no name-restriction
operator, and no ambient supply of atoms. The encoding of the
$\pi$-calculus given with the original presentation was meant to
demonstrate that $\nu$ can be discharged under that constraint --- that
one need not posit a runtime source of fresh names. Lybech showed that
the encoding is wrong, by two counterexamples, and gave a corrected
encoding that is provably valid against Gorla-style criteria. His
correction obtains freshness from a single central \emph{name server}
contacted by every $\nu$ in the term. That is a correct encoding of a
different thing: it discharges $\nu$ by supposing precisely the resource
whose elimination was the point.

We repair the encoding without a name server. The repair is not an
invention. The original construction was modelled on Abramsky's
computational interpretation of classical linear logic, and used only
half of it: the multiplicative \emph{splitting} of a namespace, and not
the \emph{rejoining} performed by the Contraction term former $t \mathbin{@} u$
together with the Copy reaction rule. Abramsky's Copy carries an
invariant --- names are pairs $\langle x_0, s\rangle$ with $s$ a binary
address, and two names sharing a root have \emph{incompatible} addresses
--- which is exactly the invariant the original encoding needed and never
stated. We show why the missing half is not optional in the rho calculus:
Copy renames the contents of a box at the metalevel, whereas rho's
duplicator copies quoted code verbatim, and rho substitution does not
descend under a quote. Address arithmetic must therefore move out of the
syntax and into processes. We give the resulting encoding, in which each
replication carries its own local address server derived from its own
address, verify it against both of Lybech's counterexamples, and prove it
fully abstract with respect to a bisimilarity whose observers are
restricted to the \emph{encodings of source contexts}. We show that this
restriction is forced rather than convenient: an unrestricted rho context
separates the encodings of $P$ and $P \mid \nil$.

Throughout we use the current concrete syntax for the rho calculus ---
$\inp{y}{x}{P}$, $\out{x}{P}$, $\quo{P}$, $\drp{x}$ --- and explain why
it, rather than the corner-quote notation of the original papers, is the
notation in which the central lemma of this note can be stated at all.
\end{abstract}

% =====================================================================
\section{Introduction}\label{sec:intro}
% =====================================================================

Process calculi customarily begin by positing a countably infinite set of
atomic names, and a scoping operator $(\nu x)P$ which conjures a name
guaranteed fresh. Both are idealisations, and both are discharged
somewhere below the model: an implementation must say where names come
from and how freshness is achieved, and in a distributed setting the
answer is not obvious.

The rho calculus \cite{MeredithRadestock05} was an attempt to move that
question inside the theory. Names are quoted processes; the only binder
is input; there is no $\nu$. Freshness, when it is available, is a
consequence of structural distinctness of the processes quoted, not of an
oracle. Whether this is a viable position is exactly the question of
whether the $\pi$-calculus can be encoded in the rho calculus
\emph{without reintroducing what was removed}, and the original paper
offered an encoding to that end. The price paid was that the translation
is not a homomorphism on a signature: it compiles a whole term at a time,
because the namespace discipline is a property of the term's structure.

Lybech \cite{Lybech22} showed that the encoding is incorrect.\footnote{%
The author discussed this with Lybech in early 2023 and mentioned that
Abramsky's construction \cite{Abramsky93} was the source of the encoding's
design and, in the author's view, the source of its repair; the thread was
not taken up. The present note makes the argument in full.}
He exhibits two counterexamples, diagnoses their common cause precisely,
proposes a new encoding, and proves it valid against a list of criteria
close to Gorla's \cite{Gorla10}, augmented by a criterion of
\emph{parameter independence} which is the right addition for
parametrised translations. All of this is correct, and this note does not
dispute any of it.

What this note disputes is the repair. Lybech's encoding obtains fresh
names from a process
\[
  !N(x,z,v,s)
\]
--- a name server --- reached over a fixed channel $v$ which every
occurrence of $\nu$ in the entire term must contact. He is candid about
the cost, and closes his paper by observing that a more elaborate
encoding, in which each replication instantiates its own source of
names, ``would be closer to the intention in the encoding by Meredith and
Radestock'', but that it appears to face a \emph{scaffolding} problem: to
build a source of fresh names one must already have a source of fresh
names. He leaves it open.

The scaffolding problem has a solution, and it was in the construction the
original encoding was modelled on. Abramsky's proof expressions for
classical linear logic \cite{Abramsky93} handle exactly this: unboundedly
many copies of a proof, each needing names distinct from every other
copy's, with no fresh names manufactured at runtime at all. The mechanism
is that a name is a pair $\langle x_0, s \rangle$ of a root and a binary
\emph{address}; the Copy rule extends addresses rather than allocating
roots; and Copy's own first clause \emph{re-contracts} the copied box's
side ports, so that the parent's namespace is never disturbed. Freshness
is required only where roots are introduced --- statically, finitely
often, at translation time.

The original encoding used the first half of this. Its splitting of the
parameter pair across parallel composition,
$\enc{P_1 \mid P_2}_{n,p} = \enc{P_1}_{+n,+p} \mid \enc{P_2}_{n+,p+}$, is
precisely Abramsky's variant operation $t \mapsto t^l, t^r$, and it is
correct for precisely his reason. It did not use the second half, and
what the second half supplies is not decoration: it is the reason the
invariant survives copying.

\paragraph{Contributions.}
\begin{enumerate}[label=(\roman*),leftmargin=2.2em]
\item We identify the invariant the 2005 encoding required and never
stated (\S\ref{sec:diagnosis}), and show that the invariant Lybech
reverse-engineers from it --- that the parameters are the most recently
replicated names --- is a different and weaker property that no splitting
clause can preserve, whereas Abramsky's prefix-incompatibility invariant
is preserved by splitting by construction. On this reading the parallel
composition clause is not the defect; it is the half that was right.
\item We prove a no-go result (Proposition~\ref{prop:nogo}) explaining why
having the right invariant is still not enough: a translation function is
a finite object and can emit only finitely many quoted addresses, while
replication yields unboundedly many occurrences; and rho's duplicator
cannot make up the difference, because it copies quoted code verbatim and
rho substitution does not descend under a quote. Address arithmetic must
therefore be performed by communication. This is stated as a syntactic
side condition on translation clauses
(Definition~\ref{def:discipline}), which both of Lybech's counterexamples
violate by inspection.
\item We give an encoding satisfying that condition
(\S\ref{sec:encoding}), in which each replication node carries a local
address server derived from its own address, the scaffolding regress
terminates at a static seed consumed once, and no process is contacted by
more than one node. Whole-term compilation survives: the address
structure is still written statically and still mirrors the parse tree;
only the addresses themselves become runtime values.
\item We verify the encoding against both of Lybech's counterexamples
(\S\ref{sec:worked}), adapted to the input-guarded fragment.
\item We prove full abstraction (\S\ref{sec:fullabs}) for a
context-labelled bisimilarity whose observers are the encodings of source
contexts, and we show (\S\ref{sec:necessity}) that the restriction is
forced: an unrestricted rho context separates $\enc{P}$ from
$\enc{P \mid \nil}$ although $P \equiv P \mid \nil$. Lybech's own
separation theorem is untouched and explains why no encoding whatever can
do better.
\end{enumerate}

\paragraph{Scope.}
The source language is the asynchronous, choice-free $\pi$-calculus with
\emph{input-guarded} replication. This is not a concession extracted by
the difficulty: unguarded replication is nonsense from an implementation
standpoint --- it runs away --- and the 2005 encoding of it diverges
unconditionally, as Lybech notes. He restricts to the same fragment for
the same reason, and observes (his footnote 5) that the restriction would
not by itself have prevented either of his counterexamples. Cost
accounting is out of scope. This note is a bug fix.

% =====================================================================
\section{The rho calculus}\label{sec:rho}
% =====================================================================

We recall the calculus \cite{MeredithRadestock05}, in current notation.

\begin{definition}[Syntax]
Processes $P$ and names $x$ are given by mutual recursion:
\[
P \ ::=\ \nil \ \mid\ P_1 \mid P_2 \ \mid\ \out{x}{P} \ \mid\ \inp{y}{x}{P}
      \ \mid\ \drp{x}
\qquad\qquad
x \ ::=\ \quo{P}
\]
$\out{x}{P}$ is the \emph{lift}: it quotes $P$, thereby creating the name
$\quo{P}$, and outputs it on $x$. $\inp{y}{x}{P}$ is the input, binding
$y$ in $P$; it is the only binder. $\drp{x}$ is the \emph{drop}: it
strips the quotes from $x$ and runs the process inside.
\end{definition}

\begin{definition}[Name equivalence and structural congruence]
Name equivalence $\nequiv$ is the least equivalence on names closed under
\[
\frac{P_1 \equiv P_2}{\quo{P_1} \nequiv \quo{P_2}}
\qquad\qquad
\frac{x_1 \nequiv x_2}{\quo{(\drp{x_1})} \nequiv x_2}
\]
and structural congruence $\equiv$ is the least congruence containing
$\alpha$-equivalence and the abelian monoid laws for $\mid$ with unit
$\nil$, with names compared up to $\nequiv$ in the congruence rules.
\end{definition}

\begin{definition}[Reduction]
\[
\frac{x_1 \nequiv x_2}
     {\inp{y}{x_1}{P_1} \ \mid\ \out{x_2}{P_2} \ \to\ P_1\{\quo{P_2}/y\}}
\]
closed under $\mid$ and $\equiv$ in the usual way.
\end{definition}

Substitution is capture-avoiding and replaces all names $\nequiv$-equal to
the substituted variable. Two clauses matter for everything below:
\[
(\quo{P})\{x/y\} =
  \begin{cases} x & \text{if } \quo{P} \nequiv y\\ \quo{P} & \text{otherwise}\end{cases}
\qquad\qquad
(\drp{x})\{\quo{P}/y\} =
  \begin{cases} P & \text{if } x \nequiv y \\ \drp{x} & \text{otherwise}\end{cases}
\]
The first is the crux of this paper: \emph{substitution does not recur
into a process under a quote}. A quoted process is an atom to
substitution; it is replaced wholesale or not at all.

\begin{notation}
We write $\out{x}{\drp{z}}$ for the transmission of the \emph{name} $z$
along $x$, since $\quo{(\drp{z})} \nequiv z$; and we write $\Dup(x)$ for
the duplicator $\inp{y}{x}{(\drp{y} \mid \out{x}{\drp{y}})}$, which
satisfies $\out{x}{Q} \mid \Dup(x) \to Q \mid \out{x}{Q}$.
\end{notation}

\subsection{Why this notation}\label{sec:notation}

The 2005 papers wrote input as $x(y).P$ and quotation with corner
quotes. We use $\inp{y}{x}{P}$, $\quo{P}$, $\drp{x}$ throughout, for four
reasons, of which the last is the one that earns its place in a paper
about a bug.

\begin{enumerate}[label=(\arabic*),leftmargin=2.2em]
\item It is ASCII. Every term displayed in this note is a term of
rholang; the appendix runs.

\item It extends to joins without renegotiation. Simultaneous consumption
from several channels is written
$\inpp{y_1}{x_1}{y_2}{x_2}{P}$, with $\&$ joining within a group and
$;$ separating alternative groups. The corner-quote prefix notation
$x(y).P$ has no comparable extension: the prefix is a unary constructor
by construction, and polyadic or joined receipt has to be simulated. The
encoding of \S\ref{sec:encoding} does not require joins, but the
generalisation of \S\ref{sec:open} to polyadic $\pi$ does, and it is
better to have paid for the notation in advance.

\item It is the for-comprehension of Scala, Haskell's \texttt{do}, and
C\#'s LINQ. A reader who knows none of the literature on process calculi
nonetheless reads $\inp{y}{x}{P}$ correctly on the first attempt: draw a
$y$ from $x$, then do $P$. The monadic reading of input is not an analogy
here; RSpace, the store underlying the rho calculus implementations, is a
comprehension over a tuplespace.

\item Most importantly, $\quo{\cdot}$ and $\drp{\cdot}$ are visibly
\emph{operators applied to terms}, occupying positions in a term. The
central technical observation of this note
(Lemma~\ref{lem:transparency}) is a statement about which \emph{positions}
substitution can reach: it reaches the object of a lift, which is a
process position, and it does not reach beneath a quote. Corner quotes
present quotation as a change of typographic register rather than as an
operator, and in that notation the distinction between
$\out{c}{\out{n}{\nil}}$ and $\quo{\out{n}{\nil}}$ --- which is the entire
difference between the broken encoding and the repaired one --- is
typographically almost invisible. It was, in fact, invisible for
seventeen years.
\end{enumerate}

\begin{remark}
The calculus is the \textbf{r}eflective \textbf{h}igher \textbf{o}rder
calculus; ``rho'' is an acronym before it is a Greek letter, and the pun
--- that rho follows $\pi$ --- is the point of the name. We therefore
never write it with the letter $\rho$.
\end{remark}

% =====================================================================
\section{The 2005 encoding and what it was for}\label{sec:2005}
% =====================================================================

Fix an injection $\varphi$ from $\pi$-calculus names to rho names; we
suppress it. The \emph{static quoting} techniques generate, from a name
$x$, the names
\[
  \Lft(x) \ \triangleq\ \quo{(\out{x}{\nil})}
  \qquad
  \Rgt(x) \ \triangleq\ \quo{(\inp{y}{x}{\nil})}
\]
(the left and right increments, written $+x$ and $x{+}$ in the original),
and their composition. Each generates a countably infinite sequence of
names from any starting point.

The translation $\enc{-}_{n,p}$ takes two name parameters, and is defined
by
\[
\begin{array}{rcl}
\enc{\nil}_{n,p} &=& \nil \\
\enc{x(y).P}_{n,p} &=& \inp{y}{x}{\enc{P}_{n,p}} \\
\enc{\bar x \langle z\rangle}_{n,p} &=& \out{x}{\drp{z}} \\
\enc{P_1 \mid P_2}_{n,p} &=& \enc{P_1}_{\Lft(n),\Lft(p)} \mid \enc{P_2}_{\Rgt(n),\Rgt(p)} \\
\enc{(\nu z)P}_{n,p} &=& \inp{z}{p}{\enc{P}_{\Lft(n),\Lft(p)}} \mid \out{p}{\drp{n}} \\
\enc{!P}_{n,p} &=& \out{a}{\big(\inp{n}{b}{\inp{p}{c}{(\enc{P}_{n,p} \mid \Dup(a) \mid \cdots)}}\big)}
                   \mid \Dup(a) \mid \cdots
\end{array}
\]
with $a = n \cdot p$, $b = \Rgt(n)$, $c = \Rgt(p)$, and top-level
parameters $n_0 = \quo{\prod_{x \in \fn(P)} \out{x}{\nil}}$,
$p_0 = \quo{\prod_{x \in \fn(P)} \inp{y}{x}{\nil}}$, chosen so that no
name of $P$ is derivable from them.

Two features of this design should be recorded before it is criticised,
because they are what the design is \emph{for}.

\begin{enumerate}[label=(\alph*),leftmargin=2.2em]
\item \textbf{There is no runtime source of names.} $n_0$ and $p_0$ are
computed by the translation from $\fn(P)$ and are never requested from
anything. Every other name in the image is obtained from them by
operations that are, in intent, local to the node that performs them.
\item \textbf{The translation compiles a whole term at a time.} It is
parametrised by a position in the term, not merely by the term's head
constructor. This is the price of (a), and it is the reason the encoding
is not a signature homomorphism.
\end{enumerate}

The defect, stated in the way we will need it, is this: the encoding
carries an invariant on $(n,p)$ that is never written down. Lybech
recovers it as \emph{the parameters are the names most recently
replicated}, and this is a fair reconstruction of what the replication
clause is trying to maintain. We shall argue in \S\ref{sec:diagnosis}
that it is the wrong invariant, and that the encoding's own ancestry
supplies the right one.

% =====================================================================
\section{Lybech's counterexamples and his repair}\label{sec:lybech}
% =====================================================================

\subsection{The counterexamples}

\begin{example}[Splitting destroys the parameter invariant]\label{ex:ce1}
Let $Q$ be arbitrary and
\[
P_1' \ \triangleq\ !\big((\nu z)\bar u \langle z\rangle \mid Q\big)
\qquad
P_2' \ \triangleq\ (\nu z)\,!\big(\bar u \langle z\rangle \mid Q\big).
\]
$P_1'$ emits a different fresh name on $u$ at each unfolding; $P_2'$
emits the same one forever, and the context
$[\,\cdot\,] \mid u(n_1).u(n_2).(\bar n_1 \mid n_2.\bar x)$
distinguishes them. Their images do not differ: the replication context
rebinds $n$ and $p$ at each unfolding, but the parallel composition
clause has already frozen $\Lft(n)$ and $\Lft(p)$ \emph{inside quotes},
where the runtime substitution cannot reach, so the inner $\nu$ emits the
same name every time.
\end{example}

\begin{example}[Nesting $\nu$ destroys it too]\label{ex:ce2}
\[
P_1 \ \triangleq\ !(\nu z)\bar u\langle z\rangle
\qquad
P_2 \ \triangleq\ !(\nu q)(\nu z)\bar u\langle z\rangle
\]
are structurally congruent, since $q$ is unused. Their images are not
equivalent: the clause for $(\nu q)$ statically increments the parameters
before the clause for $(\nu z)$ sees them, so in $\enc{P_2}$ the emitted
names are all equal while in $\enc{P_1}$ they are all distinct.
\end{example}

Lybech's diagnosis is exact, and we adopt it verbatim: the property of
being ``the most recently replicated names'' is not preserved by the
translation of parallel composition, because the increments the
translation writes are \emph{static}, i.e.\ they place the parameter
under a quote, and substitution does not recur under quotes.

\begin{remark}\label{rem:guarded}
Both examples use unguarded replication, which our fragment excludes. The
errors survive the restriction --- Lybech records as much explicitly ---
and we restate the examples in the input-guarded fragment in
\S\ref{sec:worked}, where we check the repaired encoding against them.
\end{remark}

\subsection{The name server}

Lybech's encoding is
$\enc{P} \triangleq \enc{P}_{n,v} \mid\, !N(x,z,v,s)$ where the server
\[
  !N(x,z,v,s) \ \triangleq\
  \Dup(x) \mid \out{x}{\big(\inp{a}{z}{\inp{r}{v}{(\Dup(x) \mid \out{r}{\drp{a}} \mid \out{z}{\out{a}{\nil}})}}\big)}
  \mid \out{z}{\drp{s}}
\]
holds a seed on $z$, and answers requests arriving on the fixed channel
$v$ with successive left increments of $s$. The parameter $n$ is
threaded forward through the translation and is never re-bound, which is
what removes the second error; the first is removed because access to the
name source is by a channel $v$ that never changes and so can never be
lost. The encoding is proved valid against compositionality, substitution
invariance, operational correspondence, observational correspondence,
divergence reflection, and parameter independence. We accept all of
this.

We also note the mechanism the server uses to build the successor name:
\[
  \out{z}{\out{a}{\nil}}
\]
where $a$ is a name that was \emph{received}. This computes
$\Lft(a)$ at runtime, and it works because the object of a lift is a
process position. It is the very manoeuvre the broken clauses needed and
did not use. Lybech deploys it once, inside a server, without remarking
that it is a general discipline; \S\ref{sec:diagnosis} makes it one.

\subsection{Why this is a correct encoding of a different thing}

\begin{enumerate}[label=(\arabic*),leftmargin=2.2em]
\item \textbf{It reintroduces the resource whose elimination was the
point.} A single process, reached over a single fixed channel, is the
source of every fresh name in the term. That is the $\pi$-calculus
assumption of an ambient name supply, implemented rather than discharged.
It is not a criticism of the theorem; it is an observation about what the
theorem is about.
\item \textbf{The topology is wrong.} Lybech makes this criticism
himself, and better than we would: for a distributed system with a
topology other than a star, the translation does not yield an adequate
representation of the source, because it routes every name request
through one node. The encoding preserves the semantics of a program and
not the intuitions underlying its structure.
\item \textbf{The criteria are not the right yardstick here.} A
Gorla-style list with compositionality as a primitive measures how close a
translation is to a signature homomorphism. But behavioural faithfulness
and syntactic homomorphism are independent: Milner's encoding of the
$\lambda$-calculus into $\pi$ is behaviourally faithful and is not a
signature homomorphism. Since compiling a whole term at a time is
precisely what the rho encoding trades for locality, adopting
compositionality as a primitive criterion rules the interesting design
out by fiat. In \S\ref{sec:fullabs} we take the correctness criterion to
be a single behavioural condition and derive Lybech's criteria from it.
\end{enumerate}

Finally, his own conclusion states the alternative and the obstacle: give
each replication its own copy of the name server, ``closer to the
intention'' of the original, and confront the scaffolding problem, since
to instantiate a source of fresh names one must first have a source of
fresh names. This note answers that. The answer was available in 1993.

% =====================================================================
\section{The missing half: Abramsky's construction}\label{sec:abramsky}
% =====================================================================

Abramsky's computational interpretation of classical linear logic
\cite{Abramsky93} assigns \emph{proof expressions} $\Theta;\bar t$ to
sequent proofs and gives them an operational semantics in the style of
the chemical abstract machine. We need four features of it.

\paragraph{Names carry addresses.}
In the notational interlude preceding the reaction rules, Abramsky fixes
a bijection
\[
  \mathcal{N} \ \cong\ \mathbb{N}_0 \times \{l,r\}^*
\]
so that a name is a pair $\langle x_0, s\rangle$ of a \emph{root} and a
binary \emph{address}. For a term $t$, the variants $t^l$ and $t^r$
replace every $\langle x_0, s\rangle$ occurring in $t$ by
$\langle x_0, sl\rangle$ and $\langle x_0, sr\rangle$ respectively. The
invariant established by the proof-expression assignment and maintained by
every transition is:
\begin{quote}
for all distinct names $\langle x_0,s\rangle, \langle y_0,t\rangle$
occurring in $P$, $x_0 = y_0$ implies that $s$ and $t$ are
\emph{incompatible} --- they have no upper bound in the prefix ordering.
\end{quote}

\paragraph{Contraction is a term former.}
The Contraction rule of $\mathsf{PE}_2$ is
\[
\frac{\vdash \Theta; \Gamma,\ t{:}\,?A,\ u{:}\,?A}
     {\vdash \Theta; \Gamma,\ t \mathbin{@} u : \,?A}
\]
The joiner $t \mathbin{@} u$ is a term, not a rewrite: it is a
first-class object that can be embedded in larger terms and cut against
other terms.\footnote{The rule is printed in \cite{Abramsky93} with
$u{:}\,?B$; comparison with the intuitionistic Contraction rule and with
case (12) of the realizability proof shows this to be a typo for $?A$.}

\paragraph{Copy is the rejoin.}
The reaction rule for the exponential is
\[
  \bar x(P) \perp u \mathbin{@} v
  \ \longrightarrow\
  \bar x \perp (\bar x^{\,l} \mathbin{@} \bar x^{\,r}),\quad
  \bar x(P)^l \perp u,\quad
  \bar x(P)^r \perp v .
\]
Three things are worth separating here.
\begin{enumerate}[label=(\roman*),leftmargin=2.2em]
\item Copying is \emph{demand-driven}: the box duplicates only when a
contraction node asks for two. Abramsky says as much --- the term of type
$?A^\perp$ specifies how many copies of the $!A$ term are required --- and
records that Copy is the only rule that increases the size of a proof
expression.
\item Distinctness of the two copies is by \emph{address extension}, not
by allocation. No fresh root is created. Freshness is required only where
roots are introduced, and roots are introduced only by the name
convention at the Axiom, With and Of Course rules --- statically, finitely
often.
\item The first clause, $\bar x \perp (\bar x^{\,l} \mathbin{@} \bar
x^{\,r})$, is the rejoin. The box's side ports $\bar x$ are cut against a
\emph{newly built contraction} of the two copies' corresponding ports.
The environment continues to see one port; the two copies' demands on it
are routed through a joiner. \emph{The parent's namespace is never
re-rooted.}
\end{enumerate}

\paragraph{The correspondence.}
\begin{center}
\begin{tabular}{@{}lll@{}}
\toprule
$\mathsf{PE}_2$ & 2005 encoding & status \\
\midrule
$\mathcal{N} \cong \mathbb{N}_0 \times \{l,r\}^*$ & $\Lft, \Rgt$ namespace templates & used \\
variants $t^l, t^r$ & $\enc{P_1 \mid P_2}_{n,p} = \enc{P_1}_{\Lft(n),\Lft(p)} \mid \enc{P_2}_{\Rgt(n),\Rgt(p)}$ & used \\
prefix-incompatibility invariant & --- & \textbf{not stated} \\
Contraction $t \mathbin{@} u$ & --- & \textbf{unused} \\
Copy, clause $\bar x \perp (\bar x^l \mathbin{@} \bar x^r)$ & --- & \textbf{unused} \\
demand-driven copying & eager replication (diverges) & \textbf{unused} \\
\bottomrule
\end{tabular}
\end{center}

% =====================================================================
\section{Diagnosis}\label{sec:diagnosis}
% =====================================================================

\begin{proposition}[The reconstructed invariant is the wrong one]
\label{prop:wronginv}
Let $I(n,p)$ be the property that $(n,p)$ are the names most recently
produced by the enclosing replication context. No splitting clause
preserves $I$: if $\enc{P_1 \mid P_2}_{n,p}$ delegates to sub-translations
at parameters derived from $(n,p)$, then the derived parameters do not
satisfy $I$ in the sub-terms, since they are not the most recent output of
the replication context but a function of it.
\end{proposition}

\begin{proof}
Immediate from the shape of $I$: it is a statement about a moving
frontier, and it identifies one distinguished pair globally, whereas
splitting produces two pairs neither of which is that one. It follows
that any repair which preserves $I$ must avoid deriving parameters at
$\mid$ --- which is what Lybech's threading of a single $n$ through the
translation does, at the cost of needing an external source for the
values $n$ takes.
\end{proof}

\begin{proposition}[Prefix-incompatibility is preserved by splitting]
\label{prop:rightinv}
Let $\mathrm{Inc}(S)$ be the property that the addresses occurring in a
set $S$ of names over a common root are pairwise incompatible in the
prefix ordering. If $\mathrm{Inc}(S)$ and $\sigma \in S$, then
$\mathrm{Inc}\big((S \setminus \{\sigma\}) \cup \{\sigma l, \sigma r\}\big)$.
\end{proposition}

\begin{proof}
$\sigma l$ and $\sigma r$ are incompatible with each other. For
$\tau \in S \setminus \{\sigma\}$: $\tau$ is incompatible with $\sigma$,
hence neither is a prefix of the other, hence $\tau$ is incompatible with
every extension of $\sigma$.
\end{proof}

Propositions \ref{prop:wronginv} and \ref{prop:rightinv} together say
that the parallel composition clause of the 2005 encoding is not the
defect. Against the invariant its ancestry supplies, splitting is exactly
the operation that maintains it. The defect is elsewhere.

\begin{lemma}[Substitution transparency]\label{lem:transparency}
The object of a lift is a substitution-transparent position and a quote is
not: for all $P$, $x$, $y$,
\[
  (\out{c}{P})\{x/y\} \ =\ \out{c\{x/y\}}{(P\{x/y\})}
  \qquad\text{whereas}\qquad
  (\quo{P})\{x/y\} \in \{x, \quo{P}\}.
\]
Consequently a name may be derived from a runtime-bound parameter $n$ by
the process $\out{c}{\out{n}{\nil}}$, which transmits $\Lft(n)$ on $c$;
and it may not be derived by writing $\quo{(\out{n}{\nil})}$ in the
translation.
\end{lemma}

\begin{proof}
By the definition of substitution: the lift's object is a process, and
substitution is defined by structural recursion through process
constructors; quotation is not a process constructor but the name
constructor, and its clause is the two-case clause above.
\end{proof}

\begin{corollary}[One quote level per rendezvous]\label{cor:onequote}
A name at quote depth $k$ above a runtime-received name requires at least
$k$ communications to construct.
\end{corollary}

\begin{proposition}[No static address assignment]\label{prop:nogo}
There is no translation function $\enc{-}$ of the 2005 shape --- finitely
many clauses, each emitting finitely many quoted names --- for which the
address carried by a name in the $k$-th unfolding of a replication is
written by the translation, for all $k$.
\end{proposition}

\begin{proof}
The translation is a finite object, so the set of quoted names appearing
syntactically in the image of a fixed term is finite. A replication
yields unboundedly many occurrences of its body, each of which must, if
the body contains $\nu$, contribute a name distinct from all the others.
The pigeonhole principle applies. The only remaining hope is that
substitution, applied at each unfolding, differentiates copies of a fixed
piece of syntax --- but by Lemma~\ref{lem:transparency} substitution
cannot reach beneath a quote, and by construction rho's duplicator
reproduces the quoted body verbatim. This is not an accident of the
duplicator: transmission of code without modification is what reflection
\emph{is}, and it is the same fact on which Lybech's separation theorem
turns.
\end{proof}

Proposition~\ref{prop:nogo} is the precise sense in which Abramsky's Copy
rule does not transfer. Copy renames the contents of a box at the
metalevel; rho has no metalevel operation on quoted code. Address
arithmetic must therefore leave the syntax and become communication, and
the structure that routes it is Copy's other clause: the shared port,
re-contracted per copy.

\begin{definition}[The discipline]\label{def:discipline}
A translation clause is \emph{admissible} if every parameter that will be
bound at runtime occurs in it only
\begin{enumerate}[label=(\alph*),leftmargin=2.2em,nosep]
\item in subject position of an input or a lift, or
\item within the object of a lift,
\end{enumerate}
and never within a quote written by the translation.
\end{definition}

\begin{proposition}
The clauses for $\mid$ and for $\nu$ in the 2005 encoding are
inadmissible; those for $\nil$, input and output are admissible.
\end{proposition}

\begin{proof}
By inspection: $\Lft(n) = \quo{(\out{n}{\nil})}$ places the parameter
under a quote written by the translation, and both clauses use it. The
remaining clauses do not mention the parameters except by passing them
along.
\end{proof}

This is the whole of the repair, and it is checkable in a minute. Both of
Lybech's counterexamples exploit exactly the two inadmissible clauses.

% =====================================================================
\section{The encoding}\label{sec:encoding}
% =====================================================================

\subsection{Addresses}

\begin{definition}[Address templates]
For a name $\sigma$ put
$\Lft(\sigma) = \quo{(\out{\sigma}{\nil})}$ and
$\Rgt(\sigma) = \quo{(\inp{y}{\sigma}{\nil})}$.
An \emph{address} is a name of the form $\theta(n_0)$ for
$\theta \in \{\Lft,\Rgt\}^*$ and $n_0$ the static seed.
\end{definition}

We never write $\Lft$ or $\Rgt$ in the image of the translation; they are
metalanguage, used only to \emph{name} the addresses that the encoding
produces by communication.

\begin{lemma}[Injectivity and disjointness]\label{lem:templates}
$\Lft$ and $\Rgt$ are injective on names up to $\nequiv$, their images are
disjoint, and $\QD(\Lft(\sigma)) = \QD(\Rgt(\sigma)) = 1 + \QD(\sigma)$,
where $\QD$ is Lybech's quote depth. Consequently
$\theta \mapsto \theta(n_0)$ is injective on $\{\Lft,\Rgt\}^*$.
\end{lemma}

\begin{proof}
Injectivity: $\quo{(\out{\sigma}{\nil})} \nequiv \quo{(\out{\tau}{\nil})}$
requires $\out{\sigma}{\nil} \equiv \out{\tau}{\nil}$, hence
$\sigma \nequiv \tau$; similarly for $\Rgt$. Disjointness: a lift and an
input are distinct constructors and structural congruence does not relate
them. Neither template's output is of the form $\quo{(\drp{x})}$, so
[N-DROP] does not apply and quote depth strictly increases. Injectivity of
$\theta \mapsto \theta(n_0)$ follows by induction on $|\theta|$, using
disjointness at each step and stratification to prevent an address from
coinciding with a proper prefix of itself.
\end{proof}

Lemma~\ref{lem:templates} is Abramsky's bijection
$\mathcal{N} \cong \mathbb{N}_0 \times \{l,r\}^*$, realised inside the rho
calculus with $n_0$ as the unique root. Prefix-incompatibility becomes,
here, the conjunction of two things already in the literature:
stratification by quote depth separates addresses of different length,
and template disjointness separates siblings.

\subsection{The one primitive}

Everything is built from a single macro, which obtains two fresh addresses
from a given one by communication.

\begin{construction}[$\Frsh$]
\[
  \Frsh(\sigma;\, u,\, v).\,P
  \ \triangleq\
  \inp{u}{\sigma}{\inp{v}{\sigma}{P}}
  \ \mid\ \out{\sigma}{\out{\sigma}{\nil}}
  \ \mid\ \out{\sigma}{\inp{y}{\sigma}{\nil}}
\]
binding $u$ and $v$ in $P$.
\end{construction}

$\Frsh$ is admissible: $\sigma$ occurs in subject position and in the
object of a lift, never under a written quote. After two administrative
communications, $\{u,v\} = \{\Lft(\sigma), \Rgt(\sigma)\}$, with the
assignment chosen nondeterministically. That the assignment is
nondeterministic is not a defect but the content of the design: only
distinctness is required, never identity, which is exactly Lybech's
criterion of parameter independence read as a design principle rather than
as a proof obligation.

\subsection{The clauses}

The encoding maps a $\pi$ term to an \emph{address abstraction}, a
function from a name to a rho process. We write $\enc{P}_\sigma$ for its
value at $\sigma$.

\begin{definition}[The translation]\label{def:translation}
\begin{align*}
\enc{\nil}_\sigma \ &=\ \nil \\[2pt]
\enc{\bar x\langle z\rangle}_\sigma \ &=\ \out{x}{\drp{z}} \\[2pt]
\enc{x(y).P}_\sigma \ &=\ \inp{y}{x}{\enc{P}_\sigma} \\[2pt]
\enc{P_1 \mid P_2}_\sigma \ &=\
  \inp{m}{\sigma}{\enc{P_1}_m}
  \ \mid\ \inp{m}{\sigma}{\enc{P_2}_m}
  \ \mid\ \out{\sigma}{\out{\sigma}{\nil}}
  \ \mid\ \out{\sigma}{\inp{y}{\sigma}{\nil}} \\[2pt]
\enc{(\nu z)P}_\sigma \ &=\ \Frsh(\sigma;\, z,\, m).\ \enc{P}_m \\[2pt]
\enc{!\,x(y).P}_\sigma \ &=\
  \Frsh(\sigma;\, a,\, b).\ \Frsh(b;\, s,\, c).\
  \Big(\ \Dup(a)
  \ \mid\ \out{a}{B}
  \ \mid\ \out{s}{\drp{c}}\ \Big)
\\
\text{where}\quad
B \ &=\ \inp{y}{x}{\Big(\Dup(a) \ \mid\
        \inp{c'}{s}{\ \Frsh(c';\, d,\, e).\ \big(\out{s}{\drp{e}} \mid \enc{P}_d\big)}\Big)}
\end{align*}
Top level: $\enc{P} \triangleq \enc{P}_{n_0}$ with
$n_0 = \quo{\prod_{x \in \fn(P)} \out{x}{\nil}}$.
\end{definition}

\paragraph{Reading the replication clause.}
The node splits its own address twice, obtaining a box channel $a$, a
local server port $s$, and a first seed $c$. The box $B$ is posted on $a$
and duplicated in the usual self-reproducing manner, so that
$\Dup(a) \mid \out{a}{B} \to B \mid \out{a}{B}$ and $B$ then blocks on
$x$: unfolding is input-guarded, and there is no divergence. On each
unfolding, $B$ reinstalls a duplicator, takes the current seed $c'$ off
the local port $s$, splits $c'$ into a body address $d$ and a successor
seed $e$, posts $e$ back on $s$, and runs the body at $d$.

The crucial point is what is \emph{inside} the quote. When the box is
posted, $a$, $s$, $x$ have already been substituted, because they are
bound by inputs whose continuations include the lift, and by
Lemma~\ref{lem:transparency} substitution reaches into a lift's object.
Thereafter $B$ is copied verbatim, and it needs no renaming, because it
contains no addresses --- only the port $s$ on which it asks for one.
This is Copy's first clause, $\bar x \perp (\bar x^l \mathbin{@} \bar x^r)$:
the copies share the parent's port and the routing does the
distinguishing.

\begin{remark}[No global server]
The port $s$ is fixed and shared --- among the copies of \emph{one}
replication node, and derived from \emph{that node's} address. A term
with $k$ replication nodes has $k$ mutually unaware local servers and no
channel that all of them use. The scaffolding regress terminates at
$n_0$, which is computed by the translation from $\fn(P)$ and consumed
once. Compare Lybech's single $!N(x,z,v,s)$ on a fixed $v$.
\end{remark}

\begin{remark}[Whole-term compilation survives]
Every process displayed above is written by the translation. What became
dynamic is the \emph{value} of an address, not the structure that
computes it: the tree of $\Frsh$ nodes is static and is isomorphic to the
parse tree of the source. The encoding remains a whole-term compilation
and remains non-homomorphic, which is the design (b) of \S\ref{sec:2005}.
\end{remark}

\subsection{Address allocation is a tree}

\begin{lemma}[Locality]\label{lem:locality}
In $\enc{P}_{n_0}$, each address serves as the subject of the
administrative traffic of exactly one clause occurrence, and the addresses
allocated by distinct occurrences are distinct.
\end{lemma}

\begin{proof}
Assign to each node of the source term the address at which its clause is
invoked. The clauses for $\nil$, output and input allocate nothing and
pass their address to their continuation, so an address may be shared
along a prefix chain but is consumed as a subject by at most the unique
$\mid$, $\nu$ or $!$ node terminating that chain. Each of $\mid$, $\nu$
and $\Frsh$ consumes its address $\sigma$ as a subject and allocates
exactly $\Lft(\sigma)$ and $\Rgt(\sigma)$, which by
Lemma~\ref{lem:templates} lie strictly below $\sigma$ in the address
tree. Subtrees allocated at $\Lft(\sigma)$ and $\Rgt(\sigma)$ are
disjoint. For $!$, the successive seeds are $c_0 = c$ and
$c_{k+1} = e_k$, where $d_k, e_k$ are the two children of $c_k$; hence
the body of the $k$-th unfolding occupies the subtree below $d_k$, and
these subtrees are pairwise disjoint and disjoint from every $c_j$. Note
in particular that the successor seed is a \emph{sibling} of the body
address, not an extension of it; taking the successor seed to be
$\Lft(c_k)$ while the body ran at $c_k$ would make the next seed collide
with the body's own first child.
\end{proof}

\begin{corollary}[Distinctness]\label{cor:distinct}
All names generated by $\enc{P}_{n_0}$ are pairwise non-$\nequiv$, and
none is $\nequiv$ to a name of $\fn(P)$ or derivable from one.
\end{corollary}

\begin{proof}
Lemmas~\ref{lem:templates} and \ref{lem:locality}; the last clause by the
choice of $n_0$, whose quote depth exceeds that of every $x \in \fn(P)$.
\end{proof}

\begin{lemma}[Administrative reductions terminate]\label{lem:admin}
Call a reduction \emph{administrative} if its subject is an address.
Every maximal administrative reduction sequence from $\enc{P}_\sigma$ is
finite, and no administrative reduction is observable on
$\varphi(\fn(P))$.
\end{lemma}

\begin{proof}
Each $\mid$ node contributes two, each $\nu$ node two, each $!$ node four
plus, per unfolding, one on $a$ and three on $s$ and $c_k$. Unfoldings
are guarded by an input on the source channel $x$, so no unfolding occurs
without a corresponding source communication. Hence the number of
administrative steps enabled at any point is bounded by a function of the
number of source communications performed. Unobservability is
Corollary~\ref{cor:distinct}.
\end{proof}

\begin{corollary}[Divergence reflection]
$\enc{P} \to^{\omega}$ implies $P \to^{\omega}$.
\end{corollary}

% =====================================================================
\section{The counterexamples, worked}\label{sec:worked}
% =====================================================================

We restate Examples~\ref{ex:ce1} and \ref{ex:ce2} in the input-guarded
fragment, guarding each replication with an input on a trigger channel
$w$, and supplying triggers from the context.

\begin{example}[Example~\ref{ex:ce1}, guarded]
\[
P_1' = !\,w(t).\big((\nu z)\bar u\langle z\rangle \mid Q\big)
\qquad
P_2' = (\nu z)\,!\,w(t).\big(\bar u\langle z\rangle \mid Q\big)
\]
distinguished in $\pi$ by
$C = [\,\cdot\,] \mid \bar w \mid \bar w \mid u(n_1).u(n_2).(\bar n_1 \mid n_2.\bar x)$.

In the image, the $k$-th unfolding of $\enc{P_1'}$ runs its body at
address $d_k$, and the body's parallel composition allocates
$\Lft(d_k), \Rgt(d_k)$ \emph{by communication on $d_k$}, whose value is
the $k$-th seed's child and therefore differs with $k$. The $\nu$ inside
then allocates from an address that differs with $k$, so a different name
is emitted on $u$ at each unfolding. In $\enc{P_2'}$ the $\nu$ is outside
the box; its $\Frsh$ fires once, and every unfolding emits the same $z$.
$\enc{C}$ separates them. This is the failure of Example~\ref{ex:ce1}
removed: the offending static increment has become a communication, and
the value it computes tracks the unfolding.
\end{example}

\begin{example}[Example~\ref{ex:ce2}, guarded]
\[
P_1 = !\,w(t).(\nu z)\bar u\langle z\rangle
\qquad
P_2 = !\,w(t).(\nu q)(\nu z)\bar u\langle z\rangle,
\qquad P_1 \equiv P_2 .
\]
In $\enc{P_1}$, unfolding $k$ runs
$\Frsh(d_k; z, m).\out{u}{\drp{z}}$, emitting one of
$\Lft(d_k), \Rgt(d_k)$. In $\enc{P_2}$, unfolding $k$ runs
$\Frsh(d_k; q, m).\Frsh(m; z, m').\out{u}{\drp{z}}$, emitting a child of
$m$, itself a child of $d_k$. In both cases the emitted names are
pairwise distinct across $k$, since the $d_k$ are, and in neither case is
any emitted name observable except as a $\pi$ name. The two are therefore
identified, as they must be. In the 2005 encoding they were separated
because the extra $\nu$ performed a static increment that the replication
context could not update; here it performs a communication that it can.
\end{example}

% =====================================================================
\section{Full abstraction relative to encoded contexts}\label{sec:fullabs}
% =====================================================================

\subsection{Context-labelled bisimilarity}

\begin{definition}[Context-labelled transitions]
Let $L$ be a calculus with reduction $\to$ and a class $\Ctx$ of
one-holed contexts. For $C \in \Ctx$ write
$P \xrightarrow{\,C\,} P'$ when $C[P] \to P'$, and
$P \xLongrightarrow{\,C\,} P'$ for $C[P] \to^* P'$.
\end{definition}

\begin{definition}[$\Obs$-relative bisimilarity]\label{def:dbisim}
Let $\Obs \subseteq \Ctx$. A symmetric relation $\mathcal{R}$ is a
\emph{$\Obs$-bisimulation} if $P \mathrel{\mathcal{R}} Q$ implies: for all
$C \in \Obs$, if $P \xrightarrow{\,C\,} P'$ then
$Q \xLongrightarrow{\,C\,} Q'$ for some $Q'$ with
$P' \mathrel{\mathcal{R}} Q'$. Write $\bisim_\Obs$ for the largest such
relation.
\end{definition}

On the source we take $\Obs = \Ctx_\pi$, all $\pi$ contexts, and write
$\bisim_\pi$. On the target we take $\Obs = \enc{\Ctx_\pi}$, the
\emph{encodings of source contexts}, and write $\bisim_{\enc{\Ctx_\pi}}$.

\begin{remark}[On minimality]\label{rem:minimality}
The labels here are enabling contexts, not \emph{minimal} enabling
contexts. Minimality --- relative pushouts in the sense of Leifer and
Milner, and their refinement to a subcategory of admissible observers ---
is the right long-run discipline and is developed elsewhere; nothing in
the theorem below needs it. Without minimality the label set is larger
than necessary and $\bisim_\Obs$ is correspondingly finer, which makes the
completeness direction harder rather than easier. We prove it anyway and
flag the refinement as future work.
\end{remark}

\begin{definition}[Encoding a context]
$\enc{-}$ extends to contexts by
$\enc{[\,\cdot\,]}_\sigma = [\,\cdot\,]_\sigma$, a hole expecting an
address, and by the clauses of Definition~\ref{def:translation} otherwise.
Plugging is instantiation of the address abstraction:
$\enc{C[P]} = \enc{C}\big[\enc{P}\big]$.
\end{definition}

\begin{lemma}[Compositionality]\label{lem:compositional}
$\enc{C[P]} = \enc{C}[\enc{P}]$ for all $\pi$ contexts $C$ and processes
$P$, on the nose.
\end{lemma}

\begin{proof}
By induction on $C$. Each clause of Definition~\ref{def:translation}
mentions its sub-translations only under an address abstraction, and the
hole occurs in exactly one of them.
\end{proof}

Note that this is Lybech's criterion 1 (compositionality) as a
consequence, with the coordinating context being the $\Frsh$ traffic and
its free names contained in $\varphi(\fn(P)) \cup \{\sigma\}$.

\subsection{Address independence}

\begin{lemma}[Address independence]\label{lem:addrind}
Let $\sigma, \tau$ be addresses satisfying
Corollary~\ref{cor:distinct} for $P$. Then there is a bijection $\beta$
between the names generated by $\enc{P}_\sigma$ and those generated by
$\enc{P}_\tau$, identity on $\varphi(\fn(P))$, such that
$\beta(\enc{P}_\sigma) = \enc{P}_\tau$; and consequently
$\enc{P}_\sigma \bisim_{\enc{\Ctx_\pi}} \enc{P}_\tau$.
\end{lemma}

\begin{proof}
The first claim is the substitution $\theta(\sigma) \mapsto \theta(\tau)$,
which is well defined and bijective by Lemma~\ref{lem:templates} and
total on generated names by Lemma~\ref{lem:locality}. For the second:
$\beta$ is not the identity, so the two terms are not equal; but an
encoded context uses a name it has received only as the subject of an
input or a lift, or as the object of one, and never applies $\drp{\cdot}$
to it or matches on its structure. Hence the only discriminations an
encoded observer can make between two received names are equality tests
by attempted synchronisation, which $\beta$ preserves and reflects, being
a bijection.
\end{proof}

Lemma~\ref{lem:addrind} is where the restriction to encoded observers
first does work, and it should be read carefully: the statement is
\emph{false} for unrestricted rho observers, as \S\ref{sec:necessity}
shows. This is the rho-calculus form of the $\pi$-calculus device of
bisimulation up to a bijection on generated names.

\subsection{Operational correspondence}

\begin{proposition}[Completeness]\label{prop:complete}
If $P \to P'$ then $\enc{P}_\sigma \to^{+} T$ with
$T \bisim_{\enc{\Ctx_\pi}} \enc{P'}_{\sigma'}$ for some address $\sigma'$.
\end{proposition}

\begin{proof}[Proof sketch]
By induction on the derivation of $P \to P'$. For $[\pi\text{-COM}]$,
$x(y).P_1 \mid \bar x\langle z\rangle \to P_1\{z/y\}$: the images are
$\inp{y}{x}{\enc{P_1}_{\sigma_1}}$ and $\out{x}{\drp{z}}$ at sibling
addresses, and one rho communication yields
$\enc{P_1}_{\sigma_1}\{z/y\}$, which equals
$\enc{P_1\{z/y\}}_{\sigma_1}$ because $y$ occurs in the image only in
subject and object positions (this is the substitution proposition, and it
is where admissibility is used a second time). Structural congruence:
the interesting case is $P \mid \nil \equiv P$, for which the images
differ by two administrative steps and are related by
Lemma~\ref{lem:admin} and Lemma~\ref{lem:addrind}. Replication: one
unfolding of $!x(y).P$ is matched by the box's copy step, the seed
exchange, and one $\Frsh$, all administrative, followed by the
communication on $x$. Restriction: $\Frsh$, two administrative steps.
The address bookkeeping is discharged by Lemma~\ref{lem:addrind}
throughout, which is what allows $\sigma'$ to be an arbitrary address
satisfying Corollary~\ref{cor:distinct}.
\end{proof}

\begin{proposition}[Soundness]\label{prop:sound}
If $\enc{P}_\sigma \to^* T$ then $P \to^* P'$ for some $P'$ with
$T \bisim_{\enc{\Ctx_\pi}} \enc{P'}_{\sigma'}$.
\end{proposition}

\begin{proof}[Proof sketch]
By induction on the length of the reduction, classifying each step as
administrative or as the image of a source communication
(Lemma~\ref{lem:admin} makes the classification exhaustive, since every
rho redex in the image has as subject either an address or a name of
$\varphi(\fn(P))$ or a $\nu$-generated name, and the latter two are
exactly the source channels). Administrative steps are absorbed into
$\bisim$: they are finite in number and produce no observable on
$\varphi(\fn(P))$. The only non-confluence among administrative steps is
the $\Frsh$ race, which permutes the two allocated addresses; the two
outcomes are related by Lemma~\ref{lem:addrind}.
\end{proof}

\subsection{The theorem}

\begin{theorem}[Full abstraction]\label{thm:fullabs}
For all $\pi$ processes $P, Q$,
\[
  P \bisim_\pi Q
  \qquad\Longleftrightarrow\qquad
  \enc{P} \bisim_{\enc{\Ctx_\pi}} \enc{Q}.
\]
\end{theorem}

\begin{proof}[Proof sketch]
($\Rightarrow$) Let $\mathcal{R} = \{(\enc{P}, \enc{Q}) : P \bisim_\pi Q\}$
closed under $\bisim_{\enc{\Ctx_\pi}}$ on both sides. Suppose
$\enc{P} \xrightarrow{\enc{C}} T$. If the step is administrative, $T$ is
related to $\enc{P}$ itself by Lemma~\ref{lem:admin}, and $\enc{Q}$ matches
by doing nothing. Otherwise, by Lemma~\ref{lem:compositional},
$\enc{C}[\enc{P}] = \enc{C[P]}$, and by Proposition~\ref{prop:sound} the
step is the image of $C[P] \to P_1$; since $P \bisim_\pi Q$ we have
$C[Q] \to^* Q_1$ with $P_1 \bisim_\pi Q_1$, and
Proposition~\ref{prop:complete} supplies the matching target moves.

($\Leftarrow$) Contrapositively, let $C$ separate $P$ from $Q$ in $\pi$.
Then $\enc{C}$ is a label in $\enc{\Ctx_\pi}$, and by
Lemma~\ref{lem:compositional} and Proposition~\ref{prop:complete} the
separating source computation is reproduced in the image; the observable
distinguishing $C[P]$ from $C[Q]$ is a barb on a name of
$\varphi(\fn(P) \cup \fn(Q))$, and barbs on such names are preserved and
reflected by Lemma~\ref{lem:admin} and
Corollary~\ref{cor:distinct}. Hence $\enc{P} \not\bisim_{\enc{\Ctx_\pi}} \enc{Q}$.
\end{proof}

\subsection{The restriction is forced}\label{sec:necessity}

\begin{proposition}[Unrestricted observers are too strong]
\label{prop:necessity}
There are $\pi$ processes $P, Q$ with $P \equiv Q$ --- hence
$P \bisim_\pi Q$ --- and a rho context $C$, not in $\enc{\Ctx_\pi}$, with
$C[\enc{P}] \not\bisim C[\enc{Q}]$.
\end{proposition}

\begin{proof}
Take $P = \bar u\langle v\rangle$ and $Q = \bar u\langle v\rangle \mid \nil$.
Then $\enc{P}_{n_0} = \out{u}{\drp{v}}$, which has no barb on $n_0$,
whereas
\[
\enc{Q}_{n_0} =
  \inp{m}{n_0}{\out{u}{\drp{v}}} \ \mid\ \inp{m}{n_0}{\nil}
  \ \mid\ \out{n_0}{\out{n_0}{\nil}} \ \mid\ \out{n_0}{\inp{y}{n_0}{\nil}}
\]
has a barb on $n_0$. Since $n_0$ is computed from $\fn(P)$ by the
translation, an arbitrary rho context can name it; take
$C = [\,\cdot\,] \mid \inp{k}{n_0}{\out{\mathit{succ}}{\nil}}$. No
encoded context can, because $n_0 \notin \varphi(\fn(P))$ by construction
--- which is precisely the content of the restriction.
\end{proof}

\begin{remark}[A sharper form of the same phenomenon]
The example above exploits the encoding's scaffolding. A second, deeper
form exploits reflection itself: an arbitrary rho context may apply
$\drp{\cdot}$ to a name received on $u$, running the process the name
quotes. Since a $\nu$-generated name in our encoding is $\Lft(\sigma)$ or
$\Rgt(\sigma)$ for an internal address $\sigma$, dropping it executes
$\out{\sigma}{\nil}$ or $\inp{y}{\sigma}{\nil}$ and exposes a barb on an
internal address. No encoded context does this: an encoded context uses a
received name only as a $\pi$ name, i.e.\ in subject or object position.
This is not a weakness of the present encoding. By Lybech's separation
theorem, the rho calculus manufactures new \emph{free} --- hence
observable and substitutable --- names at runtime, and no encoding of
$\pi$ into rho, his or ours, can be fully abstract against unrestricted
rho observers. The restriction to $\enc{\Ctx_\pi}$ is the shape of the
theorem, not a weakening of it.
\end{remark}

\subsection{Comparison with the $\mathcal{N}$-restricted criteria}

\begin{proposition}\label{prop:criteria}
Theorem~\ref{thm:fullabs} implies Lybech's criteria 1--6 for
$\enc{-}$, with $\simeq$ taken to be $\bisim_{\enc{\Ctx_\pi}}$ and the
name derivation function $\delta$ taken to be
$\sigma \mapsto \{\Lft(\sigma), \Rgt(\sigma)\}$.
\end{proposition}

\begin{proof}[Proof sketch]
Compositionality is Lemma~\ref{lem:compositional}; substitution
invariance is the substitution proposition used in
Proposition~\ref{prop:complete}; operational correspondence is
Propositions~\ref{prop:complete} and \ref{prop:sound}; observational
correspondence follows from Corollary~\ref{cor:distinct} and
Lemma~\ref{lem:admin}, since the observable names are exactly
$\varphi(\fn(P))$ and administrative traffic is confined to addresses;
divergence reflection is Corollary~\ref{cor:distinct}'s companion; and
parameter independence is Lemma~\ref{lem:addrind}.
\end{proof}

\begin{remark}
Lybech's $\bisim^{\mathcal{N}}$ restricts observation by \emph{stipulating}
a set of names. Ours restricts it by \emph{deriving} the observers from
the source: $\Obs = \enc{\Ctx_\pi}$. The two agree here, because $\pi$ and
rho are both name-passing calculi and barbs are an adequate notion of
observation for both. They do not agree in general, and a presentation in
terms of a name set assumes that what one can observe of a term is the
set of names it can communicate on --- true for these two calculi and
false for, e.g., the ambient calculus.
\end{remark}

% =====================================================================
\section{Comparison and discussion}\label{sec:discussion}
% =====================================================================

\begin{center}
\begin{tabular}{@{}p{0.30\textwidth}p{0.19\textwidth}p{0.21\textwidth}p{0.21\textwidth}@{}}
\toprule
 & MR 2005 & Lybech 2022 & this note \\
\midrule
runtime name source & none & one, global & none \\
topology of name traffic & --- & star & tree, $\cong$ parse tree \\
whole-term compilation & yes & no & yes \\
replication & eager, diverges & input-guarded & input-guarded \\
stated invariant & none & most-recent params & prefix-incompatibility \\
correctness criterion & full abstraction (false) & six criteria & full abstraction rel.\ $\enc{\Ctx_\pi}$ \\
status & refuted & proved & this paper \\
\bottomrule
\end{tabular}
\end{center}

Three remarks on what has and has not been achieved.

\paragraph{Freshness is confined, not eliminated.}
Abramsky's construction is not free of a freshness assumption either: his
name convention stipulates that a different root is introduced at each
instance of the Axiom, With and Of Course rules. What the construction
buys --- and what this note inherits --- is that freshness is required
only \emph{statically}, finitely often, at translation time. At runtime
there is no allocation: only address extension, performed locally, by
communication. That is the defensible form of the claim that the rho
calculus needs no theory of names from outside, and it is strictly what
the name server costs.

\paragraph{Demand-driven copying is a dividend, not a patch.}
Abramsky's Copy fires only when a contraction asks for a second copy.
Restricting to input-guarded replication --- which both Lybech and this
note do, to avoid divergence --- is the process-calculus shadow of that
discipline. It is worth recording that the fix for divergence and the fix
for freshness come from the same place.

\paragraph{Linearity was doing work that we must now do by hand.}
In $\mathsf{PE}_2$ the address invariant is protected by linearity (each
name occurs exactly twice) and acyclicity (which Abramsky uses to
guarantee deadlock freedom). The rho calculus has neither: names are used
by arbitrarily many senders and receivers. What was structural there is a
theorem here, and its proof is Lemmas~\ref{lem:templates} and
\ref{lem:locality}, which decompose along exactly Abramsky's two axes ---
depth and siblings.

% =====================================================================
\section{Related work}\label{sec:related}
% =====================================================================

Abramsky's programme was carried forward by Bellin and Scott
\cite{BellinScott94}, who translate classical linear logic proof
structures into a synchronous $\pi$-calculus, and tightened into a
session-typed discipline by Caires and Pfenning \cite{CairesPfenning10}
and Wadler \cite{Wadler14}. Kokke, Montesi and Peressotti \cite{HCP}
observe that the separation of environments internalised by $\otimes$
corresponds to parallel process behaviour, and extend linear logic to
hyperenvironments accordingly; that is the same move as the present one,
made from the side of the logic rather than the side of the encoding.

On encodability, Gorla \cite{Gorla10} supplies the criteria Lybech adapts,
and Gorla and Nestmann \cite{GorlaNestmann14} argue against full
abstraction as a criterion; van Glabbeek \cite{Glabbeek18} derives
validity from a semantic equivalence rather than a list, which is closer
in spirit to what we do here, though he does not treat parametrised
translations. Sangiorgi's encoding of HO$\pi$ into $\pi$
\cite{Sangiorgi93} is the result Lybech's separation theorem shows does
not extend to reflection. Bendixen, Bojesen, H\"uttel and Lybech
\cite{Psi} instantiate the rho calculus as a higher-order $\Psi$-calculus.
Leifer and Milner \cite{LeiferMilner00} supply the derived-transition
machinery that a minimal-label refinement of \S\ref{sec:fullabs} would
use.

% =====================================================================
\section{Open problems}\label{sec:open}
% =====================================================================

\begin{enumerate}[label=\arabic*.,leftmargin=2.2em]
\item \textbf{Minimal labels.} Replace enabling contexts by minimal
ones (relative pushouts), taken relative to the observer class
$\enc{\Ctx_\pi}$ rather than to the whole of rho, and show that
Theorem~\ref{thm:fullabs} survives with a coarser relation.
\item \textbf{Full proofs.} Propositions~\ref{prop:complete},
\ref{prop:sound}, \ref{prop:criteria} and Theorem~\ref{thm:fullabs} are
given as sketches. The bookkeeping is routine but long, and the case
that deserves scrutiny is the interaction of the $\Frsh$ race with
replication.
\item \textbf{Polyadic $\pi$ and joins.} The encoding is monadic. Polyadic
synchronisation should be encoded using the join form
$\inpp{y_1}{x_1}{y_2}{x_2}{P}$ directly, which would put
\S\ref{sec:notation}(2) to work and connect with Carbone and Maffeis
\cite{CarboneMaffeis03}.
\item \textbf{Choice.} The source fragment is choice-free.
\item \textbf{Cost.} The $\Frsh$ tree is a metering surface: each source
$\nu$ costs two communications, each unfolding four. What the encoding
charges is a natural object, and it is not treated here.
\item \textbf{Mechanisation.} Definition~\ref{def:discipline} is a
syntactic side condition on translation clauses and should be machine
checked; a translation that satisfies it cannot have a bug of the kind
found in 2022.
\end{enumerate}

\section*{Acknowledgements}

The author thanks Stian Lybech, whose counterexamples are correct and
whose diagnosis of their common cause is the reason the present repair
could be stated so briefly. Claude (Anthropic) assisted with source
review, drafting, and proof-checking.

\begin{thebibliography}{99}

\bibitem{Abramsky93} S. Abramsky.
\emph{Computational interpretations of linear logic.}
Theoretical Computer Science 111(1--2):3--57, 1993.

\bibitem{BellinScott94} G. Bellin and P. J. Scott.
\emph{On the $\pi$-calculus and linear logic.}
Theoretical Computer Science 135(1):11--65, 1994.

\bibitem{CairesPfenning10} L. Caires and F. Pfenning.
\emph{Session types as intuitionistic linear propositions.}
CONCUR 2010, LNCS 6269, 222--236.

\bibitem{CarboneMaffeis03} M. Carbone and S. Maffeis.
\emph{On the expressive power of polyadic synchronisation in $\pi$-calculus.}
Nordic Journal of Computing 10(2):70--98, 2003.

\bibitem{Glabbeek18} R. van Glabbeek.
\emph{A theory of encodings and expressiveness.}
FoSSaCS 2018, 183--202.

\bibitem{Gorla10} D. Gorla.
\emph{Towards a unified approach to encodability and separation results
for process calculi.}
Information and Computation 208(9):1031--1053, 2010.

\bibitem{GorlaNestmann14} D. Gorla and U. Nestmann.
\emph{Full abstraction for expressiveness: history, myths and facts.}
Mathematical Structures in Computer Science 26:639--654, 2014.

\bibitem{HCP} W. Kokke, F. Montesi and M. Peressotti.
\emph{Better late than never: a fully abstract semantics for classical
processes.}
Proc. ACM Program. Lang. 3(POPL), 2019.

\bibitem{LeiferMilner00} J. J. Leifer and R. Milner.
\emph{Deriving bisimulation congruences for reactive systems.}
CONCUR 2000, LNCS 1877, 243--258.

\bibitem{Lybech22} S. Lybech.
\emph{Encodability and separation for a reflective higher-order calculus.}
EXPRESS/SOS 2022, EPTCS 368, 95--112.
Technical report: Reykjav\'ik University, 2022.

\bibitem{MeredithRadestock05} L. G. Meredith and M. Radestock.
\emph{A reflective higher-order calculus.}
Electronic Notes in Theoretical Computer Science 141(5):49--67, 2005.

\bibitem{Psi} A. R. Bendixen, B. B. Bojesen, H. H\"uttel and S. Lybech.
\emph{A generic type system for higher-order $\Psi$-calculi.}
EXPRESS/SOS 2022, EPTCS 368, 43--59.

\bibitem{Sangiorgi93} D. Sangiorgi.
\emph{From $\pi$-calculus to higher-order $\pi$-calculus --- and back.}
TAPSOFT 1993, LNCS 668, 151--166.

\bibitem{Wadler14} P. Wadler.
\emph{Propositions as sessions.}
Journal of Functional Programming 24(2--3):384--418, 2014.

\end{thebibliography}

\appendix

% =====================================================================
\section{Status of results}
% =====================================================================

\begin{center}
\begin{tabular}{@{}llp{0.42\textwidth}@{}}
\toprule
Result & Status & Note \\
\midrule
Prop.~\ref{prop:wronginv}, \ref{prop:rightinv} & proved & elementary \\
Lemma~\ref{lem:transparency}, Cor.~\ref{cor:onequote} & proved & \\
Prop.~\ref{prop:nogo} & proved & \\
Lemma~\ref{lem:templates} & proved & uses Lybech's stratification lemma \\
Lemma~\ref{lem:locality}, Cor.~\ref{cor:distinct} & proved & \\
Lemma~\ref{lem:admin} & proved & \\
Lemma~\ref{lem:compositional} & proved & \\
Lemma~\ref{lem:addrind} & proved & the $\Obs$-relativity is essential \\
Prop.~\ref{prop:complete}, \ref{prop:sound} & \textbf{sketch} & bookkeeping routine; see open problem 2 \\
Thm.~\ref{thm:fullabs} & \textbf{sketch} & modulo the above \\
Prop.~\ref{prop:necessity} & proved & \\
Prop.~\ref{prop:criteria} & \textbf{sketch} & modulo operational correspondence \\
\bottomrule
\end{tabular}
\end{center}

% =====================================================================
\section{The encoding in rholang}
% =====================================================================

\begin{lstlisting}
// Fresh2(sigma; u, v).P  --- two addresses from one, by communication.
// L(sigma) = @{ sigma!(Nil) }   R(sigma) = @{ for (y <- sigma) { Nil } }

// Fresh2(sigma; u, v).P  expands inline to:
//
//   for (u <- sigma) { for (v <- sigma) { P } }
//   | sigma!( sigma!(Nil) )
//   | sigma!( for (y <- sigma) { Nil } )
//
// After two communications, {u,v} = { L(sigma), R(sigma) }.

// [[ P1 | P2 ]]_sigma
//   for (m <- sigma) { [[P1]]_m } | for (m <- sigma) { [[P2]]_m }
//   | sigma!( sigma!(Nil) ) | sigma!( for (y <- sigma) { Nil } )

// [[ (nu z) P ]]_sigma
//   for (z <- sigma) { for (m <- sigma) { [[P]]_m } }
//   | sigma!( sigma!(Nil) ) | sigma!( for (y <- sigma) { Nil } )

// [[ ! x(y).P ]]_sigma  --- local address server on s, box on a
//   Fresh2(sigma; a, b). Fresh2(b; s, c).
//     ( D(a) | a!( B ) | s!( *c ) )
//   B = for (y <- x) {
//         D(a) |
//         for (cp <- s) {
//           Fresh2(cp; d, e).( s!(*e) | [[P]]_d )
//         }
//       }
//   D(a) = for (q <- a) { *q | a!(*q) }
\end{lstlisting}

\noindent
The listing is schematic in one respect: rholang's receipts are
linear-or-persistent as a whole, and the two-message handshake of
$\Frsh$ is written above with two separate linear receives on the same
channel, which is what the calculus prescribes. A production
implementation would use the join form,
\texttt{for (u <- s \& v <- s) \{ \ldots \}}, which consumes both in one
rendezvous and removes the race of \S\ref{sec:encoding} altogether --- at
the cost of assuming joins, which the pure calculus does not have. This
is the concrete payoff of \S\ref{sec:notation}(2) and is recorded here
rather than assumed in the theory.

\end{document}
